\cotefile{157-4}
\scanpage{318}\mspage{42}
% archive sheet number: 317

\unsure{cause|comme}, que l'on a $\deleted{\unsure{A|B}}B(\mathrm{Fib}_{W})$, et a fortiori (comme par hypothèse $\unsure{\mathcal{F}|F}\subset\mathrm{Fib}_{W}$) $B(\mathcal{F})$.

\uline{NB} Il est clair en effet que si on a $\mathcal{F}\subset\mathcal{F}'$ alors
\[
 A(\mathcal{F}')\Longrightarrow A(\mathcal{F}) \qquad
 \text{et de \={m} pour }B,C. \tag*{(*)}
\]

Or la condition préliminaire sur $\mathcal{F}$ est bien satisfaite dans le cas \inserted{\unsure{d'une var.|d'un ens.} de l'\deleted{\unsure{une|forme}}des formes} $\deleted{\unsure{A}}\mathcal{F}=\mathcal{F}_{W_{\unsure{*|\bullet|0}}},\mathrm{Fibf}_{W_*},\mathrm{Parf}_{W_*}$. Il suffit \deleted{\unsure{des}} la voir dans le cas des sous-ensembles $\mathcal{F}_{W_*},\mathrm{Parf}_{W_*}$ de $\mathrm{Fibf}_{W_*}$, i.e que \deleted{\unsure{toute|tous}} \unsure{flèches|flèche} $f$ de $\mathrm{Cat}$ se factorise en
\[
 f=pi,
\]
avec $i\in W$, et $p$ étant, au choix, dans $\mathrm{Parf}_{W_*}$, ou dans $\mathcal{F}_{W_{\bullet}}$. Or dans le cas des $\mathrm{Parf}_{W_*}$, \margin{$\Vert$} c'est ce qui est établi dans $\overline{\underline{\mathrm{VII}}}$ (Catégories de chemins \unsure{(1)|(I)}), grâce au fonctorisme des $\underline{\mathrm{Ch}}(\unsure{\mathbb{S}|S}\unsure{.|;|,})\deleted{\unsure{g}}\longrightarrow\mathbb{S}\times\mathbb{S}$. Pour trouver une factorisation avec $f=qj$, $j\in W$, $q\in\mathcal{F}_{W_*}$, il suffit de voir que $p$ se factorise ainsi en $p=q\circ j_0$, on aura
\[
 f=\underset{\substack{\rotatebox{-90}{$\scriptstyle\in$}\\\mathcal{F}_{W_{\unsure{\bullet|0|*}}}}}{\underline{q\circ}}\,
 \underset{\substack{\rotatebox{-90}{$\scriptstyle\in$}\\W}}{\underline{(j_0i)}}\; ,
\]
\deleted{Mais si} avec \unsure{ops|obs} $f\in\mathrm{Parf}_{W_*}$. Mais si $X$ est $W$-parfait sur $S$, donc une
